%
% sn-konstruktion.tex -- geometrische Definition von sn, cn, dn
% (Mueller, Spezielle Funktionen 2022, Abb. 11.4)
%
% phi = am(0.9, 0.8) = 0.834564, numerisch per RK4 auf dphi/du = dn.
%
\documentclass[tikz]{standalone}
\usepackage{amsmath}
\usepackage{times}
\usepackage{txfonts}
\usetikzlibrary{calc}
%
% farben.tex -- Farben aus buch/common/macros.tex
%
\definecolor{darkgreen}{rgb}{0,0.6,0}
\definecolor{darkred}{rgb}{0.8,0,0}
\definecolor{orange}{rgb}{1,0.6,0}
\definecolor{gelb}{rgb}{1,0.8,0}

%\definecolor{darkblue}{rgb}{0,0,0.6}
\definecolor{lightblue}{rgb}{0.2,0.2,1.0}
\begin{document}
\def\skala{1.923}
\begin{tikzpicture}[>=latex, thick, scale=\skala]

% numerische Werte
\def\a{1.66667}        % grosse Halbachse 1/sqrt(1-k^2), k = 0.8
\def\cq{0.671471}      % cos(phi), phi = am(0.9, 0.8)
\def\sq{0.741031}      % sin(phi)
\def\px{1.119119}      % a*cos(phi), x-Koordinate von P

% ---------- linkes Panel: Einheitskreis ----------
\draw[->, black!45] (-1.15,0) -- (1.35,0);
\draw[->, black!45] (0,-1.35) -- (0,1.55);
\node[anchor=north east] at (1.35,0) {$x$};
\node[anchor=south east] at (0,1.33) {$y$};

\draw (0,0) circle (1);
\draw[thin] (0,0) -- (0.621610,0.783327);

% Bogen = Bogenlaenge u
\draw[darkred, line width=1.6pt]
	(1,0) arc[start angle=0, end angle=51.570, radius=1];
\node[darkred] at (25.79:1.22) {$u$};

\fill[darkred] (0.621610,0.783327) circle (0.035);
\node[darkred, anchor=south west] at (0.621610,0.783327) {$(\cos u, \sin u)$};

% ---------- rechtes Panel: Ellipse mit Hilfskreis ----------
\begin{scope}[xshift=3.4cm]

\draw[->, black!45] (-1.7,0) -- (1.9,0);
\draw[->, black!45] (0,-1.35) -- (0,1.6);
\node[anchor=north] at (1.9,0) {$x$};
\node[anchor=south east] at (0,1.38) {$y$};

\draw[dashed, thin] (0,0) circle (1);
\draw (0,0) ellipse ({\a} and 1);

% Halbachsen
\draw[thin] (0,0) -- (\a,0);
\fill (\a,0) circle (0.03);
\node[anchor=south west] at (\a,0) {$a$};
\draw[thin] (0,0) -- (0,1);
\fill (0,1) circle (0.03);
\node[anchor=south east] at (0,1) {$1$};

% Strecke vom Nullpunkt zum Fusspunkt von P
\draw[orange] (0,0) -- (\px,0);
\node[orange, below=0.1pt, inner sep=1.5pt] at ({\px/2},0)
	{$a\,\operatorname{cn}(u,k)$};

% Winkel phi
\draw[thin] (0,0) -- (\cq,\sq);
\draw[thin] (0.32,0) arc[start angle=0, end angle=47.823, radius=0.32];
\node at (17:0.40) {$\varphi$};

% sn(u,k)
\draw[darkgreen] (\px,0) -- (\px,\sq)
	node[midway, sloped, below=0.1pt, inner sep=0.5pt]
	{$\operatorname{sn}(u,k)$};

% cn(u,k)
\draw[lightblue] (0,\sq) -- (\cq,\sq)
	node[midway, above=0.1pt, inner sep=0.5pt] {$\operatorname{cn}(u,k)$};

\draw[dashed, thin] (\cq,\sq) -- (\px,\sq);

% r = a dn(u,k)
\draw[darkred] (0,0) -- (\px,\sq)
	node[pos=0.55, below=0.2pt, darkred] {$r$};
\node[darkred, anchor=south west] at ($(\px,\sq)+(-0.2,0.3)$)
	{$r = a\,\operatorname{dn}(u,k)$};

\fill (\cq,\sq) circle (0.028);
\node[anchor=north, yshift=1.5pt] at (\cq,\sq) {$Q$};

\fill[darkred] (\px,\sq) circle (0.035);
\node[darkred, anchor=south] at (\px,\sq) {$P$};

\end{scope}

\end{tikzpicture}
\end{document}
